\documentclass[../main.tex]{subfiles}
\begin{document}
%==========================================TIÊU ĐỀ TẬP========================================
\begin{table}[h]
  \begin{adjustwidth}{-1cm}{}
    \begin{tabular}{>{\centering\arraybackslash}p{6cm}|>{\raggedright\arraybackslash}p{12.5cm}}
      \multirow{5}{6cm}
      % Dòng đầu Yêu tinh: ÉPISODE và số tập
      {\\[-40pt]
        \flaregothic\fontsize{20pt}{20pt}\selectfont \centering
        \color{eptype}ÉPISODE\color{black}\\
        \burbank\fontsize{72pt}{72pt}\selectfont
        \color{epnum}91
        \\[18pt]
        \flaregothic\fontsize{14pt}{14pt}\selectfont
        \color{epattr}BIENVENUE, \uline{\color{Wamy} Lamy} !\\
        % \flaregothic\fontsize{18pt}{18pt}\selectfont
        % \color{epattr}ÉPISODE PERDU
        \color{black}
      }
       &
      % Dòng thứ nhất: tên tập tiếng Nhật
      {\fotskip\fontsize{18pt}{18pt}\selectfont \ruby{白}{しろ}\ruby{熊}{くま}\ vs. お\ruby{嬢}{じょう}\ruby{様}{さま}
      } \\[6pt]
      % Dòng thứ hai: tên tập tiếng Anh
       & {\cafeteria\fontsize{18pt}{18pt}\selectfont Polar Bear vs. Noble Girl}                         \\[4pt]
      % Dòng thứ ba: tên tập tiếng Việt
       & {\vnmsans\fontsize{18pt}{18pt}\selectfont Yêu tinh tuyết bợm rượu}                             \\[8pt]
      % Dòng thứ tư: Nhân vật xuất hiện
       & {\brian{\fontsize{14pt}{14pt}\selectfont Track used: Frostbite Caves - Ultimate Battle}
         }
      \\[8pt]
      % Dòng cuối cùng: Thời gian diễn ra của tập (thường sẽ là ngày phát hành)
       & {\continuum\fontsize{16pt}{16pt}\selectfont Ngày phát hành: 31/01/2021}                        \\[18pt]
    \end{tabular}
  \end{adjustwidth}
\end{table}
%=========================================TRUYỆN CHÍNH========================================
\color{black}

\pov{Hikawa Kuon}

Một tuần nữa đã qua đi, và thành viên tiếp theo của \textit{NePoLaBo} lại tới.

Một nàng Yêu tinh tuyết đến từ vùng đất băng giá, mang trong mình vẻ đẹp thanh khiết nhưng cũng ấm áp hơn những gì người ta tưởng. Nếu Pol-pol là một nghệ sĩ xiếc rực rỡ với những hành động không thể dự đoán được, thì thành viên này lại là một bông hoa tuyết nhẹ nhàng, luôn mang đến cảm giác bình yên cho những ai ở bên cạnh.

Em ấy có một sức hút rất riêng - dù bề ngoài trông như một tiểu thư dịu dàng, nhưng thực chất lại là một người vô cùng kiên định. Với giọng nói êm ái như gió đông và sự quan tâm chân thành, em ấy khiến người khác cảm thấy như được bao bọc trong một tách trà nóng giữa ngày đông lạnh giá.

Tính cách của em ấy cũng làm cho người ta chú ý: dịu dàng, chu đáo khi ở bên những người thân thiết, nhưng lại có chút nhút nhát khi gặp người lạ. Bề ngoài trông có vẻ mong manh, nhưng thực chất lại là một người rất chăm chỉ và luôn nỗ lực để trở thành phiên bản tốt nhất của chính mình.

Ngoài ra, em ấy còn là một chuyên gia\dots\ uống rượu sake. Đối thủ để có thể đấu tửu lượng với bố tôi - người cũng là người có nhiều chiêu trong các buổi nhậu nhẹt (nhưng ông ấy không phải là người bợm rượu) đây.

Và dù là một nàng Yêu tinh tuyết, mục tiêu của em không chỉ là mang đến sự thanh thản và êm ái như một cơn gió lạnh trong mùa đông năm nay. Em ấy muốn một ngày nào đó có thể trở thành một nghệ sĩ thực thụ, để giọng nói của mình sưởi ấm trái tim của mọi người, dù là trong những ngày đông giá rét nhất.

Một nàng Yêu tinh với giấc mơ lấp lánh.

\dots Hoặc có lẽ, một ``bông tuyết hoang dã'' sẵn sàng tan chảy vì đam mê của mình.

\chrint

Yukihana Lamy (\ruby{雪}{ゆき}\ruby{花}{はな}ラミィ, \ruby{Giu-ki}{Tuyết}-\ruby{ha-na}{Hoa}\ La-mi\footnote{Còn có thể đọc là Ra-mi.}) là thành viên thứ hai của Thế hệ thứ Năm. Con gái của một gia đình quý tộc từ một vùng đất tuyết xa xôi. Sau khi xem được những thước phim, những buổi sóng trực tiếp của \hll, em ấy quyết định rời khỏi quê hương cùng với người bạn đồng hành là Daifuku (con gấu lật đật) để gia nhập vào công ty ``đào tạo thần tượng kiêm diễn viên hài ngoài giờ'' này.

\begin{wrapfigure}[25]{l}{8.5cm}
  \includegraphics[width=8.5cm]{../../../../Image/Book?/lamy_model.png}
\end{wrapfigure}

Tính cách của Lamy-chan có thể được mô tả ngắn gọn là ``trong sáng'' và ``yêu thương tất cả,'' bởi có chất giọng nhẹ nhàng, là một người tốt bụng và dịu dàng, đồng thời thể hiện sự đồng cảm khi tương tác với mọi người - từ đồng nghiệp cho tới những người xem của em ấy, vì thế biệt danh ``Lamy-mama'' cũng từ đó mà ra.

Mặc dù lúc đầu em ấy có vẻ nhút nhát và e thẹn, nhưng rồi nhanh chóng nhận ra rằng em ấy có thể táo bạo như họ vốn có và không ngại khẳng định mình khi tình huống đòi hỏi.

Lamy-chan thân thiết với những người thuộc Thế hệ thứ Năm - nhất là Botan\footnote{Sẽ được giới thiệu ở tập sau.}. Cả hai được biết đến với những lần chơi trò chơi kinh dị, với một số người nói đùa rằng việc Botan-chan hành động như một người ``chồng'' bảo vệ người ``vợ'' Lamy-chan đang sợ hãi (mặc dù Botan rõ ràng thích phản ứng của Lamy).

Cũng như đã trình bày ở trên, dù là một người ``seiso'' (trong sáng), Lamy-chan cũng có những hành động khá là\dots\ đáng ngờ, ví dụ có những khoảnh khắc dâm đãng, đôi lúc thì thể hiện tính cách yandere, và hơn hết, là rất thích uống rượu. Bố tôi mà biết được rằng tôi có một người bạn nhậu xứng tầm với ông ấy thì chắc cả hai người sẽ uống rượu cả đêm đây.

Vì Lamy-chan là một Yêu tinh (giống như Flare-chan, mặc dù em ấy có mang nửa dòng máu là người) nên việc có đôi tai nhọn hoắt là chuyện đương nhiên. Em có đôi mắt vàng, có mái tóc xanh lam rất dài đến giữa lưng, với phần đuôi nhuộm màu xanh đậm hơn, phần mái chéo chia đôi dài đến giữa mắt và có chỏm tóc ngố tạo hình trái tim.

Kiểu tóc thường ngày của Lamy-chan bao gồm một cặp bím tóc nửa buộc nửa thả được buộc bằng ruy băng thành kiểu đuôi ngựa thấp, có đeo một bông hoa cài tóc và đồ trang trí tóc hình bông tuyết.

Bộ trang phục của em ấy bao gồm một chiếc mũ nồi trắng, một chiếc áo sơ mi trắng không tay có diềm xếp nếp ở giữa và đường cắt xẻ ở giữa ngực, một chiếc nơ họa tiết ca rô màu xanh, một chiếc áo khoác lông họa tiết bông tuyết màu xanh (thường hay mặc trễ vai), một chiếc áo nịt ngực màu đen, một chiếc thắt lưng da màu nâu đeo lỏng, một chiếc váy ngắn xếp ly màu xanh (có chút ánh hồng ở dưới) diềm xếp nếp mang. Lamy-chan đi tất dài đến đùi có diềm xếp nếp màu trắng và đôi bốt cài dây cao gót màu nâu-đen có viền lông.

\chrint

Nghe nói rằng thành viên này là ``Mẹ hiền'' của \hll, chỉ sau Mio-chan. Không biết người mẹ hiền này mang tới trò gì tới văn phòng nào\dots

\il

Ngày cuối cùng của tháng Một, thời tiết rét căm căm. Không biết có phải bởi vì mùa đông đã dần lên đến đỉnh điểm, hay bởi vì thành viên mới này đấy.

Bên ngoài, từng cơn gió lạnh lùa qua khung cửa sổ, mang theo hơi thở cuối cùng của mùa đông. Mặt trời vẫn chưa lên hẳn, bầu trời xám xịt báo hiệu một ngày lạnh lẽo và lười biếng.

Bên trong căn phòng ấm cúng, lò sưởi đang cháy tí tách, tỏa ra ánh sáng cam dịu nhẹ. Dẫu vậy, bầu không khí lại chẳng hề yên bình như vẻ ngoài của nó.

``Rừm! Rặc! Rặc! Rặc! Rặc! Rặc! Rặc!'' tiếng của Lamy-chan vang lên, phá tan cả sự im lặng vốn có.

\gif{7-9i}{1}{82}

Tôi thở dài, chống tay lên trán: ``Dừng lại đi em ơi.''

Tôi nhìn nàng Yêu tinh tuyết đang vừa rung người, vừa ngước mắt lên trần nhà, miệng không ngừng phát ra những âm thanh khó hiểu. ``Em kêu nãy giờ được nửa tiếng rồi đấy. Dừng lại đi.''

Nàng ta chẳng có vẻ gì là nghe thấy tôi nói, vẫn tiếp tục tạo những âm thanh kỳ quặc kia.

Mio-chan và Flare-chan, hai nạn nhân vô tội của tình huống này, đã lặng lẽ trốn ra sau kệ tủ, chỉ ló mỗi cái đầu ra để quan sát từ xa.

\img{7-9b}

Nàng Dị giới cau mày: ``Lamy-chan đang làm cái quái gì thế?''

Nàng Sói đen liếc đồng đội, giọng thản nhiên như thể đã quen với cảnh tượng này: ``Chắc là đang bắt chước tiếng máy cắt cỏ vì đang chán ấy mà.''

Tôi nhíu mày: ``Lúc buồn chán bộ em ấy toàn làm thế suốt à\dots?''

Flare-chan thở dài: ``Có đầy thứ mà em ấy có thể bắt chước được\dots''

Tôi và Mio-chan nhìn nhau, cùng đồng thanh chốt hạ: ``\dots mà lại bắt chước cái thứ của nợ đó. Lạ ghê.''

Sau một hồi phát ra những âm thanh vô nghĩa đến mức tôi bắt đầu tự hỏi liệu em ấy có bị ``lỗi hệ thống'' không, nàng Yêu tinh tuyết cuối cùng cũng chịu dừng lại. Cô thở phào, lấy ra một chai rượu từ đâu đó rồi tự lẩm bẩm: ``Được rồi, nghỉ một chút nào.''

Nàng Dị giới ló mặt ra, thốt lên, tay chỉ về chai rượu mà Lamy-chan đang cầm kia với vẻ mặt sửng sốt: ``Khoan đã, em lấy cái đó ở đâu ra vậy!?''

Tôi không ngạc nhiên. Sơ yếu lý lịch của Lamy đã cho tôi biết trước rằng em ấy là một ``bợm rượu'' chính hiệu. Nhưng thay vì trả lời, Lamy-chan chỉ ôm chặt lấy chai rượu như thể nó là báu vật, rồi kêu lên hoảng hốt: ``C-Chị không được uống đâu!''

Tôi và Flare-chan liếc nhau, đồng thanh đáp lại: ``Ai thèm uống của em chứ! Nhất là vào giữa ban ngày nữa!''

Nhận lại chỉ là một nụ cười trừ từ người đồng hương kia.

Tôi nhìn chằm chằm vào Lamy-chan đang ôm chặt chai rượu như thể nó là báu vật cuối cùng trên thế gian. Nhìn đi nhìn lại, tôi mới nhận ra một điều: cái chai này không phải loại bình thường. Cái nhãn ấy\dots\ không phải là của mấy loại rượu có thể mua ngoài siêu thị.

Không biết em ấy kiếm được cái chai rượu kia bằng cách nào nhỉ?

``Chị nghe nói là em ấy hay giấu rượu ở mọi ngóc ngách trong văn phòng này á,'' Mio-chan đột nhiên lên tiếng, giọng điệu như thể đang tiết lộ một bí mật động trời.

Cả ba người chúng tôi ngay lập tức quay đầu nhìn về phía ``lính mới''. Cảm giác như một con thú nhỏ vừa bị phát hiện giữa bầy thợ săn, Lamy-chan chỉ biết lùi lại từng chút một, mắt đảo qua đảo lại tìm đường thoát. Nhưng đã quá muộn.

``Thế thì bắt đầu kiểm tra xem nào.'' Tôi và Flare-chan nhìn nhau, rồi bước tới như hai thanh tra đầy kinh nghiệm. Mục tiêu đầu tiên, tất nhiên, là chai rượu trong tay của nàng Yêu tinh tuyết.

\begin{figure}[H]
  \centering
  \begin{subfigure}[t]{0.45\textwidth}
    \centering
    \includegraphics[width=8cm]{../../../../Image/Book?/7-9c1.png}
  \end{subfigure}
  \quad
  \begin{subfigure}[t]{0.45\textwidth}
    \centering
    \includegraphics[width=8cm]{../../../../Image/Book?/7-9c2.png}
  \end{subfigure}
  \quad
  \begin{subfigure}[t]{0.45\textwidth}
    \centering
    \includegraphics[width=8cm]{../../../../Image/Book?/7-9c3.png}
  \end{subfigure}    \quad
  \begin{subfigure}[t]{0.45\textwidth}
    \centering
    \includegraphics[width=8cm]{../../../../Image/Book?/7-9c4.png}
  \end{subfigure}
\end{figure}

Flare-chan giơ tay lên, chỉ vào nó:
``Cái này?''

Mio-chan gật đầu chắc nịch.
``Cái đó.''

Em ấy tiếp tục mở hộc tủ ngay gần đó, lôi ra một chai khác:
``Vậy, cái này cũng thế?''

Nàng Sói đen vẫn giữ nguyên vẻ mặt bình thản:
``Đúng.''

Tôi cúi xuống, kéo ngăn kéo ti-vi ra và giơ lên một chai nữa:
``Cả cái này nữa?''

Mio-chan không thèm suy nghĩ, đáp ngay:
``Cũng thế.''

Nàng Yêu tinh tóc vàng với tay lên bàn máy tính, nhấc một chai nữa lên, lần này bắt đầu có chút hoang mang:
``Và cái này?''

Mio vẫn gật đầu đều đều:
``Cũng thế nốt.''

Tôi liếc sang những cái bình và chai đủ loại được sắp xếp khắp phòng, chỉ tay:
``Mấy cái bình với cái chai này nữa?''

Mio-chan nhún vai như thể chuyện này quá bình thường:
``Vâng.''

Flare-chan lật đật quay qua, nhấc thanh kiếm của Ayame-chan đang đặt dựa vào tường:
``Cả cái này nữa à?''

Nàng Sói đen nghiêm túc gật đầu:
``Ừ.''

Lúc này thì chúng tôi thực sự phải nhìn nhau với vẻ nghi ngờ. ``Đến cái này mà cũng chứa rượu được à?'' chúng tôi đổ mồ hôi hột, lẩm bẩm với sự khó tin.

Bỏ qua sự tồn tại khó hiểu của một thanh kiếm có thể trở thành bình chứa rượu, Flare-chan tiến đến một chiếc hộp bìa to đặt ở góc phòng, từ từ mở nó ra, như thể đang mong chờ một bí mật kinh thiên động địa nào đó.

Nhưng thay vì rượu như chúng tôi dự đoán, bên trong lại là\dots

\dots một nàng Cừu đang ngủ ngon lành.

\img{7-9d}

Ba chúng tôi nghiêng đầu đầy khó hiểu, rồi nhìn nhau như muốn hỏi: ``Tại sao em ấy lại nằm trong này?''

Không ai biết phải nói gì. Cuối cùng, Yêu tinh tóc vàng chỉ lặng lẽ\dots\ đóng chiếc hộp lại, như thể chuyện này chưa từng xảy ra.

Nhưng ngay khoảnh khắc chúng tôi để lộ vị trí mà Watame-chan đang nằm, Mio-chan chợt sững người.

Tôi có thể cảm nhận được bầu không khí xung quanh đột nhiên thay đổi. Đôi tai Sói đen-chan giật giật, ánh mắt ánh lên một tia sáng kỳ lạ.

\dots Hiển nhiên, bản năng săn mồi nguyên thủy của em ấy lại trỗi dậy. Chậc.

\ct

Tôi nhìn Lamy-chan với ánh mắt đầy hoài nghi. Nếu tôi không nhầm, thì em ấy vừa cố gắng chuồn khỏi đây với lý do mà đến chính em ấy cũng chưa chắc đã tin nổi.

``Mà, em không còn thời gian ngồi ì ở đây nữa đâu! Em sẽ về nhà và làm một buổi stream uống rượu buổi tối!''

Cái cớ thật lộ liễu.

Tôi mỉa mai đáp lại: ``Lý do nghe hợp lý nhể.''

Flare-chan thì bùng nổ ngay lập tức. Em ấy sán lại gần, lắc đầu và chéo tay như đang khiêu khích: ``Em chỉ muốn uống thôi, đúng không? Mơ đi.''

Như bị đâm trúng tim đen, mặt nàng Lamy-chan tối sầm lại. Đôi mắt em ấy như đang muốn xuyên thủng chúng tôi, nhưng đáng tiếc, ánh nhìn đó chẳng có chút hiệu lực nào cả.

Người đàn chị khoanh tay, nhún vai đầy khiêu khích: ``Em có lườm đến mấy thì chị cũng sẽ không lui đâu!''

\begin{figure}[H]
  \centering
  \begin{subfigure}[t]{0.45\textwidth}
    \centering
    \includegraphics[width=8cm]{../../../../Image/Book?/7-9e1.png}
  \end{subfigure}
  \quad
  \begin{subfigure}[t]{0.45\textwidth}
    \centering
    \includegraphics[width=8cm]{../../../../Image/Book?/7-9e2.png}
  \end{subfigure}
\end{figure}

Mọi chuyện có lẽ sẽ còn tiếp tục nếu như chúng tôi không nhận ra một mối nguy khác.

Sói đen-chan đang từ từ tiếp cận chiếc hộp bìa ban nãy. Em ấy cúi thấp người, mắt dán chặt vào chiếc hộp. Đôi tai nhọn khẽ giật giật, cái mũi hơi động đậy, như thể ngửi thấy mùi gì đó.

Tất nhiên, nàng Dị giới cũng nhận ra ngay, liền hỏi: ``Chị đang định làm cái gì đấy, Senpai?''

Mio-chan giật mình, như thể vừa bị bắt quả tang đang làm chuyện mờ ám. Em ấy quay đầu lại, miệng lắp bắp một lý do mà chính cô ấy có lẽ cũng thấy miễn cưỡng: ``Chị đang định chuẩn bị miếng thịt này vào bữa tối.''

Cả hai chúng tôi lập tức phản ứng: ``Cái đó có ăn được đâu!''

Tôi thì trợn mắt, không thể tin nổi vào tai mình: ``Em biết em không phải là sói thuần, đúng chứ?''

Mio-chan rên rỉ như một đứa trẻ bị tước mất món đồ yêu thích: ``Nhưng mà\dots\ trông em ấy ngon quá\dots! Em muốn được ăn thịt em ấy\~{}!''

Tôi day trán, thở dài ngao ngán: ``Kìm hãm tính thú của em đi. Kẻo bị đuổi việc chứ chả chơi đấy.''

Mio chỉ biết rên rỉ đầy tiếc nuối. Mối nguy hiểm giờ đã qua đi, còn giờ\dots\ chúng tôi lại phải đối mặt với một vấn đề khác: tên bợm rượu.

Lamy-chan  đang ôm chặt chai rượu, bực bội thốt lên: ``Vậy là ba người anh chị quyết cản em buổi uống rượu tối của em ha\dots''

Tôi và Flare-chan còn đang chật vật giữ nàng Sói đen lại khỏi Watame-chan kia, nhưng vẫn phải quay qua phản ứng.

\img{7-9f}

Nàng Dị giới nhướng mày ngạc nhiên: ``Sao phải cáu thế?''

Nàng Sói đen, dù đang bị giữ lại, vẫn cố lên tiếng: ``Ể? Sao chị cũng bị lôi vào nữa?''

Tôi chỉ biết ngán ngẩm thở dài: ``Cái quái\dots? Cái này là chuyện bình thường mà? Đây có phải giờ uống rượu đâu?''

Thấy con Sói đen đã bị khống chế, Lamy-chan hiểu rằng không thể dựa vào tình huống hỗn loạn để lẩn đi được nữa. Cô nàng chống nạnh, bĩu môi tỏ vẻ bực bội, nhưng rồi ánh mắt chợt lóe lên như vừa nghĩ ra điều gì đó. Em ấy khẽ hắng giọng, cố lấy lại khí thế: ``\textbf{Ra đây nào, Daifuku!}''

\img{7-9g}

Cùng với tiếng gọi dõng dạc ấy, từ góc phòng lăn ra một con lật đật gấu trắng Bắc Cực đồ chơi. Nó lắc lư qua lại, trên cổ quấn một chiếc khăn màu xanh đặc trưng. Tôi nheo mắt quan sát, rồi chợt nhớ ra một chi tiết mà mình đã đọc trong hồ sơ: ``\textit{Hình như tôi đã nói ở phần giới thiệu rằng đây là bạn đồng hành của em ấy, đúng không nhỉ?}''

Mio-chan bật cười đầy hứng thú: ``Ô, Daifuku-chan này.''

Tôi khoanh tay, nhướng mày nhìn con gấu lật đật vừa xuất hiện kia: ``Con gấu lật đật đó thì làm được gì-''

Nhưng chưa kịp nói hết câu, nàng Yêu tinh tuyết đã giơ tay hô lớn: ``\textbf{Biến hóa!}''

Ngay lập tức, một luồng ánh sáng xanh nhạt tỏa ra từ con lật đật. Và rồi, khi ánh sáng dịu lại, trước mặt chúng tôi là một con gấu Bắc Cực to lớn\dots\ nhưng có gì đó rất sai sai.

\img{7-9h}

Flare-chan trợn mắt: ``Tại sao nó lại vác cái máy cắt cỏ chứ!?''

Tôi há hốc mồm nhìn cảnh tượng khó hiểu trước mắt, rồi vỗ trán: ``Không lẽ đây là tàn dư từ sự chán nản của em ấy?''

Mio-chan rên rỉ, đưa tay bóp trán: ``Làm ơn biến trở lại bình thường đi cái.''

Dẫu vậy, con gấu đó vẫn thản nhiên vẫy tay chào chúng tôi. Nó không có vẻ nguy hiểm gì, chỉ là\dots\ trông rất ngớ ngẩn và vô lý khi một con gấu Bắc Cực cầm theo một chiếc máy cắt cỏ như thể sẵn sàng đi làm vườn bất cứ lúc nào.

Ngay sau đó, Lamy-chan chỉ tay về phía chúng tôi và hô lớn: ``Được rồi, Daifuku! Bắt ba người bọn họ cùng với chị uống rượu nào!''

Tôi chưa kịp phản ứng thì con gấu lập tức hành động. Nhưng thay vì lao đến giữ chặt chúng tôi hay ép chúng tôi cầm cốc rượu, nó lại\dots\ đi nhặt hết mấy chai rượu mà chủ nhân của nó đã giấu khắp nơi trong phòng.

Cả ba chúng tôi đứng đơ ra như tượng, không hiểu chuyện gì đang diên ra trước mắt.

Nàng Dị giới chớp mắt liên tục: ``Khoan\dots\ chuyện quái gì đang xảy ra vậy?''

Còn chưa hết sốc, chúng tôi lại chứng kiến cảnh tượng đáng kinh ngạc hơn: con gấu Bắc Cực không chỉ gom hết rượu lại một chỗ, mà còn thản nhiên giật luôn cả chai rượu mà chính chủ của nó đang cầm trên tay và giơ lên cao như thể muốn nói: ``Không được uống!''

\gif{7-9i}{83}{175}

Lamy-chan sững sờ nhìn cảnh tượng trước mắt, rồi như bị chọc trúng dây thần kinh nào đó, em ấy lập tức nhào tới, đấm túi bụi vào người con gấu. Với lực đấm hầu như không đáng kể.

``\textbf{Waah!} Em đang làm cái gì thế, Daifuku!? Dừng lại đi!''

Nhưng con gấu vẫn đứng yên, không hề hấn gì, chỉ cúi đầu xuống đầy kính cẩn về phía chúng tôi, như muốn thay mặt chủ nhân gửi lời xin lỗi vì ``đã có một cô chủ bợm rượu tới mức thảm hại như vậy,'' ít nhất là từ những gì tôi đọc được.

Tôi ôm bụng cười ngả nghiêng, còn nàng Yêu tinh tuyết thì vừa đấm vừa khóc như một đứa trẻ bị cướp mất kẹo. Cứ tưởng em ấy sẽ dùng con gấu để phản công chúng tôi, ai ngờ chính ``trợ thủ'' của em lại trở mặt một cách đầy phũ phàng.

Flare-chan liếc nhìn cả hai với vẻ khó hiểu, rồi hỏi: ``Vậy, ờm\dots\ Daifuku thực sự là cái gì vậy?''

Tôi nhún vai, đáp lại một cách thản nhiên: ``Trong hồ sơ lý lịch thì nó là biểu trưng cho fan-club của em ấy. Nó cũng là người kiểm soát chế độ uống rượu của em ấy nữa.''

Mio-chan chớp mắt, có vẻ ngạc nhiên: ``Thế ạ? Mà thôi, không cần phải lo chuyện này nữa đâu.''

Tôi nhìn nàng Yêu tinh tuyết vẫn đang vật vã với con gấu mà thở dài: ``Thật. Trông cũng tội, mà thôi cũng\dots\ kệ vậy.''

Lamy-chan ngước lên, đôi mắt long lanh như sắp khóc thêm lần nữa, miệng bĩu ra đầy ấm ức: ``Diệu của chụy mừ\~{}!'' và tiếp tục đấm túi bụi bạn trợ thủ của mình.

\gif{7-9j}{1}{11}

Tôi bật cười. Quả nhiên, trận chiến này vẫn có người thắng, mà không ai ngờ kẻ chiến thắng lại là một con gấu lật đật.
%==========================================TRUYỆN PHỤ=========================================
\omk

``Trả cho chị đây! Trả đây\~{}!! \textbf{Daifuku!!!}''

Lamy-chan liên tục đấm vào thân con gấu trắng khổng lồ của mình, nhưng trông chẳng khác nào muỗi đốt inox. Daifuku vẫn đứng im như tượng đá, kiên định không chịu trả lại dù chủ nhân của nó đang giãy giụa đầy tức tối. Tôi khoanh tay nhìn cảnh tượng trước mặt, khóe miệng nhếch lên. Nhìn em ấy vừa đánh vừa khóc đòi rượu thế này, phải công nhận là vừa đáng yêu vừa buồn cười.

Mio-chan ngập ngừng, tay chống cằm: ``Chúng ta có định ngăn cản em ấy lại không\dots?''

Flare-chan nhún vai, không mấy bận tâm: ``Hay là cứ kệ em ấy vậy.''

Tôi bật cười: ``Có cái hài để xem cũng vui mà.''

Yêu tinh vàng liếc tôi với ánh mắt không mấy đồng tình: ``Em lại chả thấy buồn cười tí nào\dots''

``Hmm\dots'' đến lượt nàng Sói đen khoanh tay đầy suy tư.

Khi nhìn nàng Yêu tinh tuyết vẫn đang khổ sở đấm đá ``thú cưng'' của mình, tôi chợt nhớ đến chuyện khác cũng liên quan đến rượu. Một ý tưởng lóe lên trong đầu khiến tôi hơi nhếch mép: ``Hay là anh sẽ mời bố anh với cả mấy đứa trong lớp anh một bữa nhậu nhỉ?''

Flare quay sang tôi, nhướng mày đầy nghi hoặc: ``Thật luôn à\dots? Đến cả Kuon-senpai cũng nghĩ đến rượu sao?''

Mio-chan cũng tò mò: ``Anh định sẽ đi nhậu ở đâu?''

Tôi nhún vai như thể đó là chuyện hiển nhiên:  ``Ở nhà anh. Nếu không uống rượu được thì không cần phải đến đâu.''

Nghe vậy, nàng Dị giới nheo mắt, có vẻ hơi do dự: ``Nghe nói là Kawachi có văn hóa nhậu khá lạ\dots''

Nàng Sói đen gật gù: ``Chị cũng từng được nghe anh ấy kể về văn hóa quê anh ấy rồi. Hay là thôi, ta cứ đi thử nhỉ?''

Flare khoanh tay, thở dài một hơi rồi gật đầu: ``Chắc em chỉ uống nước ngọt với nước trái cây quá\dots''

Tôi nhún vai, chấp nhận: ``Cũng được. Coi như là một dịp để hiểu thêm nhà anh một tí. Vậy là hai đứa đến hả?''

Cả hai gật đầu xác nhận. Tôi ``hừm'' một tiếng đầy ẩn ý, rồi vỗ tay một cái, gọi lớn: ``\textbf{Ê, Daifuku!}''

Tôi nghĩ chắc cũng muốn trêu em ấy một chút nhỉ?

\pov{Yukihana Lamy}

Nghe tiếng gọi, con gấu của tôi ngoảnh mặt về phía giọng nói, còn tôi thì lập tức dừng đánh nó lại. Mặc dù vẫn còn uất ức vì bị ``phản bội'', tôi cũng muốn biết xem anh ta định làm gì.

``Đi được với bọn này chứ?'' anh ấy lớn giọng hỏi.

Tôi cau mày, chớp mắt khó hiểu. ``Đi đâu vậy anh?''

Nhưng Kuon-senpai không nhìn tôi, mà vẫn nhìn chằm chằm vào bạn đồng hành của tôi: ``Anh không có hỏi em. Anh hỏi Daifuku.''

Tôi trừng mắt về phía anh ấy, thảng thốt: ``Cái gì chứ!?''

Con gấu của tôi nhìn anh ấy với vẻ băn khoăn, như thể vẫn chưa hiểu ý. Anh ấy liền nói tiếp: ``Bọn tao đi nhậu. Mày có thể uống được không?''

Daifuku lắc đầu. Hiển nhiên rồi, con gấu của tôi làm sao có thể uống rượu được! Chỉ có Yukimin mới có thể làm được thôi!

Kuon-senpai tặc lưỡi, có vẻ hơi thất vọng. ``Đành vậy.'' Nhưng ngay sau đó, nó lại cầm một cái cốc ra, mở chai rượu mà nó đã giật của tôi khi nãy và rót vào cốc một cách rất thuần thục. Tôi và anh ấy đều tròn mắt nhìn.

Daifuku, tại sao em lại làm thế với chị chứ!? Đó là chai rượu ngon nhất mà chị để dành đấy!

Anh ấy nghiêng đầu hỏi lại: ``Tức là mày có thể rót rượu cho bọn tao chứ gì?''

Con gấu của tôi gật đầu. Không lẽ Kuon-senpai định rủ bạn đồng hành của tôi đi cùng sao?

Anh ấy bật cười. ``À, vậy thì phiền mày nha!'' Quả nhiên là như vậy. Chị sẽ không chịu để em đi như vậy đâu, Daifuku!

Tôi lập tức lên tiếng: ``Khoan đã! Chị cũng muốn đi nữa! Kuon-senpai! Em cũng muốn đi--''

Chưa để tôi nói hết câu, bạn đồng hành của tôi\dots\ lắc đầu. Nó giơ tay ra hiệu ``không,'' đồng thời chặn đứng tôi trước mặt.

Tôi đứng sững tại chỗ, sốc đến mức không nói nên lời khi bị chính bạn đồng hành phản bội lần nữa. ``\textbf{\dots Tại sao chị lại không được đi chứ!?}''

Kuon-senpai khoanh tay, cười nửa miệng đầy trêu chọc. ``Anh nghĩ là em là kẻ bợm rượu đó, Lamyssian-chan.''

Tôi giật mình. ``Em không có bợm rượu! Mà `Lamyssian' là thế nào chứ!?''

Anh ấy nhún vai. ``Thì dân Nga người ta cũng uống rượu nhiều như em đấy.''

Flare-senpai đứng bên cạnh khẽ gật đầu. ``\dots Cái đó em phải công nhận anh nói đúng, Kuon-senpai ạ.''

Mio-senpai cũng cười khẽ: ``Chẳng biết anh ấy học chơi chữ ở đâu nữa\dots''

Tôi nghiến răng, lườm Daifuku, rồi đấm nhẹ vào thân nó một cái nữa: ``Chị là chủ nhân của em mà, sao em dám ngăn chị uống rượu hả!?''

Anh ấy phì cười với câu mắng của tôi: ``Thôi, anh nghĩ là em nên về nhà đi stream uống rượu một mình đi. Chính em lúc nãy còn mở miệng đòi về mà. Các Yukimin đang chờ em đấy.''

Tôi siết chặt tay, dậm chân giãy nảy, lăn ra ăn vạ: ``Không chịu đâu\~{}!! Muốn đi, muốn đi, muốn đi cơ!''

Kuon-senpai bật cười trước phản ứng của tôi, còn hai người tiền bối kia thì chỉ lắc đầu ngán ngẩm, miệng lẩm bẩm: ``Đúng là kẻ bợm rượu\dots''

Cuối cùng, có lẽ vì thấy tôi quá thảm thương, anh ấy cũng đành xuống nước: ``Thôi được rồi, đi thì đi. Nhưng đừng có làm loạn ở nhà anh là được.''

Nghe vậy, tôi lập tức tươi tỉnh, nhoẻn miệng cười: ``Vậy mới đúng chứ!''

Daifuku thở dài, lắc đầu ngán ngẩm. Tôi giả vờ như không nhìn thấy.

\ct

Đêm hôm đó, tại nhà Kuon-senpai, không khí vô cùng náo nhiệt. Ngoài tôi, Flare-senpai và Mio-senpai, anh ấy còn rủ thêm cả vài người bạn nữa - mà thú thực, tôi cũng chẳng nhớ nổi tên họ.

Bữa tiệc rượu bắt đầu trong sự hào hứng, mọi người cười nói rôm rả, nâng ly liên tục. Ban đầu, tôi còn giữ ý một chút, nhưng rồi không khí vui vẻ khiến tôi cũng chẳng còn giữ kẽ làm gì. Tôi và bố của anh Tổng - Ryujin-san - bắt đầu so tài uống rượu.

Giờ phút này, xung quanh tôi là một ``bãi chiến trường'' toàn người gục ngã. Flare-senpai thì ngồi tựa vào bàn, đôi mắt lờ đờ, lẩm bẩm gì đó không rõ. Mio-senpai thì đang ngủ gật, còn mấy người bạn của Kuon-senpai thì chẳng ai còn tỉnh táo. Cả Daifuku của tôi - giờ trở lại làm con thú dễ thương như bình thường - cũng đã ngủ trong lòng tôi từ khi nào.

Tôi ngồi đối diện với Ryujin-san, chậm rãi đặt chén rượu xuống.

Tôi cười đắc thắng, nhìn người đàn ông vẫn còn khá tỉnh táo: ``Hể\~{} Vẫn chưa gục à, Ryujin-san?''

Ryujin-san nhướng mày, cười nhẹ: ``Cô bé, cháu cũng khá lắm đấy. Bố không nghĩ con có thể tìm được một đối thủ xứng tầm đâu, Kuon à.''

Ông ấy rót thêm rượu vào chén của tôi, rồi chạm ly với tôi một cách dứt khoát. Trong lúc tôi và bố anh ấy vẫn còn vui vẻ bên nhau, Kuon-senpai dường như cũng đã bắt đầu lim dim đôi mắt.

Anh ấy chống cằm, hỏi chúng tôi với giọng yếu ớt: ``Hai người\dots\ còn trụ nổi thật à?''

Tôi vỗ ngực tự hào đáp lại: ``Tất nhiên! Đẳng cấp của một Yêu tinh tuyết mà!''

Ryujin-san bật cười: ``Cháu không tệ, nhưng chưa thắng được ta đâu, cô bé.''

Tôi nheo mắt nhìn ông ấy, như thể chấp nhận lời thách đấu. ``Vậy thì tiếp tục nào, Ryujin-san!''

``Được, ta cũng muốn xem giới hạn của cháu đến đâu!'' nói đoạn, chúng tôi cụng ly thêm một cái nữa, rồi uống hết một hơi.

Kuon-senpai thở dài, lẩm bẩm: ``Hai người này đúng là y như nhau\dots'' rồi gục xuống bàn, để lại tôi và Ryujin-san tiếp tục màn đấu tửu lượng không hồi kết.

\il

Cuối cùng, dù tôi đã uống rất sung sức, nhưng vẫn không thể đánh bại được Ryujin-san. Tôi không hiểu nổi tại sao một người trông có tuổi như vậy lại có thể giữ được tửu lượng đáng kinh ngạc. Khi tôi bắt đầu cảm thấy chóng mặt và gục xuống bàn, ông ấy chỉ cười nhẹ, rồi chậm rãi đặt ly xuống như thể chẳng hề hấn gì.

Sáng hôm sau, khi tỉnh dậy, tôi nhận ra mình đang nằm trong một căn phòng xa lạ - à không, là nhà của Hikawa-san. Có vẻ như tôi đã uống quá đà đến mức ngủ lại luôn ở đây. Lúc ngồi dậy, đầu tôi vẫn hơi ong ong, nhưng dù vậy, tôi vẫn không thể quên được trận đấu tửu lượng hôm qua.

Mấy ngày sau, khi tôi đến văn phòng \hll, tôi vẫn còn nhớ rất rõ về bữa nhậu đó. Đúng là hai cha con nhà này chẳng giống người bình thường chút nào\dots

Tôi chợt nhớ lại hình ảnh Kuon-senpai gục trên bàn, rồi ông bố của anh ấy vẫn ung dung rót thêm một ly rượu nữa. Hai người bọn họ tuy khác nhau về phong cách, nhưng đều khiến tôi cảm thấy có một sự\dots\ đáng tin cậy theo cách riêng của họ.

Tôi còn chưa kịp nghĩ ngợi thêm thì cửa văn phòng chợt mở ra. Một giọng nói vang lên, kéo tôi về thực tại.

``Chào mọi người\~{} Tôi đến rồi đây!''

Tôi quay đầu lại, và ngay lập tức nhận ra người sắp gia nhập vào ``bãi chiến trường'' văn phòng này vào tuần sau. Một nụ cười nở trên môi tôi - hẳn là nơi này sắp trở nên náo nhiệt hơn rồi.

Một cô nàng Sư tử trắng. Có ngoại hình khá ngầu lòi. Trùm bắn tỉa trong các trò FPS.

Kuon-senpai, hy vọng là anh sẽ sẵn sàng tiếp đón thành viên tiếp theo trong nhóm \textit{NePoLaBo} chúng em\dots\ Sẽ đáng mong chờ đây.

\end{document}